\section{Step 21 Continue through the complete byte interval}\label{crc-s21}

Advance one byte, reduce the remaining length, and return to the table guard. The whole update now shares the bit routine interface. It can receive a raw state produced by either method and must preserve the same caller memory and saved-register contract.

\Needspace{11\baselineskip}
\subsection{Write and inspect}
Replace the marked temporary instructions in \texttt{t\_next}. Retain surrounding labels and markers.

\MARSCodeLabel{MIPS assembly}
\begin{minted}[samepage=true]{mars-idea}
    daddiu $a0, $a0, 1
    addiu $a1, $a1, -1
    beq $zero, $zero, table_next_byte
    nop
\end{minted}

Use \texttt{s21-1}, \texttt{s21-2}. Prepare the exact memory inputs in a separate source file and pause at \texttt{kernel\_done}, before the caller stores the result.

\subsection{Expected observation}
The default table update returns raw \texttt{340bc6d9}, with no source or table stores.

\Needspace{11\baselineskip}
\subsection{Check your understanding}
May final correctness require the illustrated temporaries?

\begin{itemize}
\item \textbf{A} No; the public state, memory and ABI contract is what matters.
\item \textbf{B} Yes, because \texttt{t0} is the checksum format.
\item \textbf{C} Only for a zero checksum.
\end{itemize}
